\documentclass[12pt]{article}
\usepackage[T1]{fontenc}
\usepackage[utf8]{inputenc}
\usepackage{lmodern}
\usepackage{textcomp}
\usepackage{enumitem}
\usepackage{geometry}
\geometry{margin=1in}
\begin{document}
\begin{center}
Answer Key - Economics - Graduate Level Assessment 18 \
\small (Answers are ordered exactly as in the question paper.)
\end{center}

\section*{Multiple Choice}
\begin{enumerate}[label=\arabic*.]
\item \textbf{Answer:} A
\\ \textit{Options:} A) Seeking to increase the demand for the union's labor; B) Negotiations to obtain a wage floor; C) Encouraging workers to increase their productivity; D) Lobbying for higher minimum wage laws; E) Promoting worker training and education
\\ \textit{Note:} The correct answer is that seeking to increase the demand for the union's labor is not a method used by unions to increase wages, as unions typically focus on negotiating higher wages and improving working conditions rather than directly influencing labor demand, which is often determined by broader market forces.
\item \textbf{Answer:} D
\\ \textit{Options:} A) Domestic steel producers lose revenue.; B) Allocative efficiency is improved.; C) It results in a decrease in domestic steel production.; D) Income is transferred from domestic steel consumers to domestic steel producers.; E) Domestic steel prices fall.
\\ \textit{Note:} A protective tariff on imported steel raises the price of foreign steel, benefiting domestic steel producers through increased sales and revenue while simultaneously imposing higher costs on domestic consumers who must pay more for steel products.
\item \textbf{Answer:} A
\\ \textit{Options:} A) natural unemployment.; B) underemployment.; C) discouraged worker.; D) structural unemployment.; E) jobless recovery.
\\ \textit{Note:} This is an example of natural unemployment because it reflects the seasonal nature of certain jobs, where employment is lost due to predictable changes in the economy, such as the closure of a seasonal facility like a public swimming pool.
\item \textbf{Answer:} A
\\ \textit{Options:} A) changes in the money supply have significant effects.; B) fiscal policy should not be used to manage the economy.; C) the FED should allow the money supply to grow at a constant rate.; D) fiscal policy has no impact on the economy.; E) the FED should decrease the money supply to combat inflation.
\\ \textit{Note:} According to Keynesian theory, changes in the money supply can influence interest rates and aggregate demand, thereby significantly affecting overall economic activity and output.
\item \textbf{Answer:} C
\\ \textit{Options:} A) Government intervenes only during economic crises to stabilize the market; B) Government plays a minor role by enforcing contracts and property rights; C) No government intervention in the economy; D) Government provides public goods and services but does not control private enterprise; E) Limited government intervention primarily in defense and public safety
\\ \textit{Note:} In a pure capitalist system, the key concept is that the economy is driven by free markets and private ownership, which inherently minimizes or eliminates the role of government intervention in economic activities.
\end{enumerate}

\section*{True / False}
\begin{enumerate}[label=\arabic*.]
\setcounter{enumi}{5}
\item \textbf{Answer:} False
\\ \textit{Options:} A) True; B) False
\\ \textit{Note:} Resource and product prices are primarily determined by market forces of supply and demand, rather than solely by government regulations and policies.
\item \textbf{Answer:} False
\\ \textit{Options:} A) True; B) False
\\ \textit{Note:} The statement is false because the United States would not export rice when the world price is lower than the domestic price, as exporters would only seek to sell in the global market when they can receive a higher price than what is available domestically.
\item \textbf{Answer:} True
\\ \textit{Options:} A) True; B) False
\\ \textit{Note:} Regulating a monopoly to produce at the allocatively efficient quantity aligns the quantity produced with consumer demand, thereby reducing the deadweight loss that occurs when the monopoly restricts output below this level to maximize profits.
\item \textbf{Answer:} False
\\ \textit{Options:} A) True; B) False
\\ \textit{Note:} The statement is false because, with a 10 percent reserve ratio, the maximum money creation from a $500 deposit can be calculated using the money multiplier effect, which allows for total money creation of up to $5,000.
\item \textbf{Answer:} False
\\ \textit{Options:} A) True; B) False
\\ \textit{Note:} In a perfectly competitive market, marginal revenue remains constant and equal to the market price for each additional unit sold, rather than increasing as output increases.
\end{enumerate}

\section*{Long Form Response}
\begin{enumerate}[label=\arabic*.]
\setcounter{enumi}{10}
\item \textbf{Answer:} A capital/output ratio of 2 implies that for every unit of output produced, two units of capital are utilized, indicating a capital-intensive production process. This suggests that the annual output will be relatively low in comparison to the level of capital investment, potentially leading to higher fixed costs and lower flexibility in adjusting to market changes. When this ratio is adjusted for quarterly output, one would expect the capital/output ratio to remain consistent, assuming that the capital stock does not change and production processes operate efficiently. However, if production is ramped up or down in response to demand fluctuations, the quarterly output could reflect a more variable capital utilization rate, leading to a potential decrease in the capital/output ratio if output increases more significantly than capital investment within that timeframe. Overall, understanding this ratio is crucial for evaluating production efficiency and investment strategies in an economy.
\\ \textit{Note:} correct because it accurately describes the implications of a capital/output ratio of 2 as indicative of a capital-intensive production process, which results in lower annual output relative to capital investment and suggests that this ratio would remain consistent when adjusted for quarterly output, assuming no changes in production methods or market conditions.
\item \textbf{Answer:} In economic terms, "land" encompasses not only traditional natural resources like minerals and soil but also renewable resources, specifically wind and solar energy. This broader definition recognizes the increasing importance of sustainable energy sources in the modern economy. By classifying wind and solar energy as forms of land, economists highlight their role as vital inputs in production processes, which can be harnessed without depleting the resource base. This inclusion underscores the necessity of integrating renewable resources into economic planning and resource management, reflecting the shift towards sustainability in economic theory and practice. Thus, the economic definition of land captures the full spectrum of natural resources essential for contemporary growth and environmental stewardship.
\\ \textit{Note:} correct because it expands the economic definition of "land" to include renewable resources like wind and solar energy, emphasizing their significance as sustainable inputs in production that do not deplete the resource base.
\end{enumerate}

\vfill
\noindent\textit{End of Answer Key}
\end{document}